% !TeX program = xelatex
%% =====================================================================
%%  ใบกิจกรรมในชั้นเรียน (Worksheet) บทที่ 1 : ตรรกศาสตร์ --- สัปดาห์ที่ 2
%%  สำหรับแจกนักศึกษาทำในห้อง (ไม่มีเฉลย มีช่องว่างให้เขียน)
%% =====================================================================
\documentclass[a4paper,12pt]{article}
\usepackage{fontspec}
\usepackage{amsmath,amssymb}
\usepackage[margin=2cm]{geometry}
\usepackage{enumitem}
\usepackage{array}
\usepackage{xcolor}

\XeTeXlinebreaklocale "th_TH"
\XeTeXlinebreakskip = 0pt plus 1pt
\setmainfont[Script=Thai,Scale=1.25]{TH Sarabun New}
\definecolor{navy}{HTML}{1F3A93}
\renewcommand{\arraystretch}{1.5}
\setlist{topsep=4pt,itemsep=6pt}

%% เส้นว่างให้เขียนคำตอบ (#1 = จำนวนบรรทัด)
\newcounter{bl}
\newcommand{\blank}[1]{\par\vspace{0.2em}%
  \setcounter{bl}{0}\loop\ifnum\value{bl}<#1\relax
  \noindent\textcolor{gray!60}{\hrulefill}\par\vspace{1.0em}\stepcounter{bl}\repeat}

\begin{document}
\thispagestyle{empty}

%% ---------- หัวกระดาษ ----------
\begin{center}
  {\Large\bfseries\color{navy} ใบกิจกรรมในชั้นเรียน บทที่ 1 : ตรรกศาสตร์ (สัปดาห์ที่ 2)}\\[0.2em]
  {\normalsize รายวิชา 4091606 คณิตศาสตร์สำหรับคอมพิวเตอร์}
\end{center}
\vspace{0.3em}
\noindent\begin{tabular}{@{}p{0.5\textwidth} p{0.22\textwidth} p{0.22\textwidth}@{}}
  ชื่อ-สกุล: \dotfill & รหัส: \dotfill & กลุ่ม: \dotfill
\end{tabular}
\hrule
\vspace{0.6em}
\noindent\textbf{คำชี้แจง:} ให้นักศึกษาทำกิจกรรมต่อไปนี้ลงในใบกิจกรรม \textbf{ภายในเวลาที่กำหนด}
แสดงวิธีคิดให้ชัดเจน แล้วร่วมเฉลยและอภิปรายพร้อมกันในชั้นเรียน
\vspace{0.6em}

%% ---------- กิจกรรม 1 ----------
\noindent\textbf{กิจกรรมที่ 1} \hfill \textit{(6 นาที)}\\
จงจำแนกว่าแต่ละประพจน์เป็น \textbf{สัจนิรันดร์ / ข้อขัดแย้ง / ประพจน์กลาง}
\begin{enumerate}[label=\arabic*)]
  \item $p \lor \neg p$ \hfill \dotfill
  \item $p \land \neg p$ \hfill \dotfill
  \item $(p \rightarrow q) \rightarrow (\neg q \rightarrow \neg p)$ \hfill \dotfill
  \item $p \rightarrow q$ \hfill \dotfill
\end{enumerate}
\vspace{0.4em}

%% ---------- กิจกรรม 2 ----------
\noindent\textbf{กิจกรรมที่ 2} \hfill \textit{(6 นาที)}\\
จงเขียนนิเสธให้อยู่ในรูปที่ ``นิเสธอยู่ติดตัวแปร'' โดยใช้กฎเดอมอร์แกน
\begin{enumerate}[label=\arabic*)]
  \item $\neg(p \land \neg q)$ \hfill \dotfill
  \item $\neg(\neg p \lor q)$ \hfill \dotfill
  \item $\neg(p \lor (q \land r))$ \hfill \dotfill
\end{enumerate}
\vspace{0.4em}

%% ---------- กิจกรรม 3 ----------
\noindent\textbf{กิจกรรมที่ 3} \hfill \textit{(8 นาที)}\\
จงเติมตารางค่าความจริงเพื่อตรวจว่า $p \rightarrow (q \land r)$ สมมูลกับ $(p\rightarrow q)\land(p\rightarrow r)$ หรือไม่
\begin{center}\small
\begin{tabular}{|c|c|c|c|c|c|c|c|}
  \hline
  $p$ & $q$ & $r$ & $q\land r$ & $p\rightarrow(q\land r)$ & $p\rightarrow q$ & $p\rightarrow r$ & $(p\rightarrow q)\land(p\rightarrow r)$ \\ \hline
  T & T & T & & & & & \\ \hline
  T & T & F & & & & & \\ \hline
  T & F & T & & & & & \\ \hline
  T & F & F & & & & & \\ \hline
  F & T & T & & & & & \\ \hline
  F & T & F & & & & & \\ \hline
  F & F & T & & & & & \\ \hline
  F & F & F & & & & & \\ \hline
\end{tabular}
\end{center}
สรุป: \dotfill
\vspace{0.4em}
\end{document}
